\begin{scanpage}{0039}{22}
where
\[v_1^{\pm}(\phi,\phi_0)=\frac i{2\pi}\sin\phi_0\int_{D_0}\frac{e^{ikr\cos\zeta}\,d\zeta}{\cos\phi_0+\cos(\zeta\pm\phi)}.\tag{3.20}\]
It is not difficult to see that with $\alpha-\pi<\phi_0<\pi$, when both faces of the wedge are illuminated, the field excited by the second face $(\phi=\alpha)$ will be described by the same equations if one replaces $\phi$ by $\alpha-\phi$, and $\phi_0$ by $\alpha-\phi_0$ in them.

Adding the field being radiated by the uniform part of the current with the incident plane wave (1.23), we obtain the diffraction field in the physical optics approach. It equals
\[\frac{E_z}{E_{0z}}=\begin{cases}v_1^+(\phi,\phi_0)+e^{-ikr\cos(\phi-\phi_0)}-e^{-ikr\cos(\phi+\phi_0)}\\\hfill\text{with }0<\phi<\pi-\phi_0,\\v_1^+(\phi,\phi_0)+e^{-ikr\cos(\phi-\phi_0)}\\\hfill\text{with }\pi-\phi_0<\phi<\pi,\\v_1^-(\phi,\phi_0)+e^{-ikr\cos(\phi-\phi_0)}\\\hfill\text{with }\pi<\phi<\pi+\phi_0,\\v_1^-(\phi,\phi_0)\hfill\text{with }\pi+\phi_0<\phi<\alpha\end{cases}\tag{3.21}\]
if one face of the wedge $(0<\phi_0<\alpha-\pi)$ is illuminated, and
\[\frac{E_z}{E_{0z}}=\begin{cases}v_1^+(\phi,\phi_0)+v_1^-(\alpha-\phi,\alpha-\phi_0)+\\\qquad+e^{-ikr\cos(\phi-\phi_0)}-e^{-ikr\cos(\phi+\phi_0)}\\\hfill\text{with }0<\phi<\pi-\phi_0,\\v_1^+(\phi,\phi_0)+v_1^-(\alpha-\phi,\alpha-\phi_0)+\\\qquad+e^{-ikr\cos(\phi-\phi_0)}\quad\text{with }\pi-\phi_0<\phi<\alpha-\pi,\\v_1^+(\phi,\phi_0)+v_1^+(\alpha-\phi,\alpha-\phi_0)+\\\qquad+e^{-ikr\cos(\phi-\phi_0)}\quad\text{with }\alpha-\pi<\phi<\pi,\\v_1^-(\phi,\phi_0)+v_1^+(\alpha-\phi,\alpha-\phi_0)+\\\qquad+e^{-ikr\cos(\phi-\phi_0)}\quad\text{with }\pi<\phi<2\alpha-\pi-\phi_0,\\v_1^-(\phi,\phi_0)+v_1^+(\alpha-\phi,\alpha-\phi_0)+\\\qquad+e^{-ikr\cos(\phi-\phi_0)}-e^{-ikr\cos(2\alpha-\phi-\phi_0)}\\\hfill\text{with }2\alpha-\pi-\phi_0<\phi<\alpha,\end{cases}\tag{3.22}\]
\end{scanpage}
